\documentclass[12 pt]{article}

\usepackage{stoversymb}
\usepackage[margin=0.75in, top=0.875in, bottom=0.875in]{geometry}
\usepackage{amsmath,amsfonts,amssymb,url,multicol,graphicx,tikz,soul}
%===makes urls render well===
\usepackage{lmodern}
\usepackage[T1]{fontenc}
%============================
\usepackage{wasysym} % smileys
\usepackage[inline]{enumitem}

%\everymath{\displaystyle}

\makeatletter
\renewcommand*\env@matrix[1][\arraystretch]{%
	\edef\arraystretch{#1}%
	\hskip -\arraycolsep
	\let\@ifnextchar\new@ifnextchar
	\array{*\c@MaxMatrixCols c}}
\makeatother

%\setenumerate{itemsep=0.25in}
%\setlist[enumerate,1]{label=\arabic*., itemsep=0.5in}
%\setlist[enumerate,2]{label={(\alph*)}, itemsep=0.8in}
%\setlist[enumerate,2]{leftmargin=0.5in}
\setlist[enumerate,3]{label={\roman*.}}
\setlist[itemize,1]{label=$\circ$, itemsep=0.25in}

\newcommand{\truefalse}[1]{#1\hfill\rule[-1mm]{220pt}{0.75pt}}
\newcommand{\hint}[1]{\vspace{3mm}\textbf{Hint}: #1}
\newcommand{\note}[1]{\textbf{Note}: #1}
\newcommand{\infsum}[3]{\sum_{{#1}={#2}}^\infty {#3}}
\newcommand{\axes}[1]{\begin{center}\includegraphics[scale=#1]{3DAxes}\end{center}}
\newcommand{\comps}[1]{\langle #1_1,#1_2,#1_3\rangle}
\newcommand{\compslong}[3]{\langle #1, #2, #3\rangle}
\newcommand{\ijk}[2]{#1\vect{#2}}
\newcommand{\pmat}[1]{\begin{pmatrix}[1.25]#1\end{pmatrix}}
\newcommand{\bmat}[1]{\begin{bmatrix}[1.25]#1\end{bmatrix}}

\setlength\parskip{0.75mm}	

\graphicspath{ {./../img/} }
\DeclareGraphicsExtensions{.pdf}

\begin{document}
%\begin{flushright}Name: \line(1,0){200}\end{flushright}
\begin{center}
\Large{\textbf{Elementary Row Operations, REF, and RREF}}
\end{center}
%\textbf{Directions:} Split into groups and (attempt to) answer each of the following questions with your group members. \ul{Don't worry about questions you are unable to answer!}

Now that we know a little bit about matrices, we're going to learn how to use matrices to solve problems!

One of the most useful things we can do to a matrix is to ``row reduce'' it. 

Row reduction is a process by which a matrix is simplified into an ``equivalent'' matrix which is easier to use overall. In order to make this procedure more canonical, we'll perform our reduction using a very precise collection of operations known as \ul{elementary row operations}.

Throughout, we'll refer to the matrix
\[
\sansM=\pmat{4 & 1 & -2 & 0 \\ 1 & 0 & 0 & -5 \\ 0 & 1 & 3 & -1 \\ 1 & 1 & -1 & -1}
\]
for all of our examples.

\subsection*{Elementary Row Operations}
There are \textbf{three} elementary row operations that we can perform on a matrix to get a new matrix which is considered ``row equivalent'' to it:\footnote{\textbf{Definition:} Two matrices $\sansM$ and $\sansN$ are said to be \ul{row equivalent} if there is a series of elementary row operations which transforms $\sansM$ into $\sansN$ (and vice versa).} 

\begin{enumerate}[itemsep=0.25in]
	\item \ul{Interchange two rows}.
	
	\textbf{Example:} Interchanging rows 2 and 4 (shorthand: $\text{R2}\leftrightarrow\text{R4}$) in $\sansM$ yields the following:
	\[
	\pmat{4 & 1 & -2 & 0 \\ 1 & 0 & 0 & -5 \\ 0 & 1 & 3 & -1 \\ 1 & 1 & -1 & -1} \xrightarrow{\text{R2}\leftrightarrow\text{R4}}
	\pmat{4 & 1 & -2 & 0 \\ 1 & 1 & -1 & -1 \\ 0 & 1 & 3 & -1 \\ 1 & 0 & 0 & -5}
	\]
	
	\item \ul{Multiply all entries in a row by a nonzero constant}. 
	
	\textbf{Example:} We can multiply row 1 of $\sansM$ by 3 (shorthand: $\text{R1}=3\text{R1}$) to get the following:
	\[
	\pmat{4 & 1 & -2 & 0 \\ 1 & 0 & 0 & -5 \\ 0 & 1 & 3 & -1 \\ 1 & 1 & -1 & -1} \xrightarrow{\text{R1}=3\text{R1}}
	\pmat{12 & 3 & -6 & 0 \\ 1 & 0 & 0 & -5 \\ 0 & 1 & 3 & -1 \\ 1 & 1 & -1 & -1}
	\]
	
	\newpage
	
	\item \ul{Add (or subtract) a nonzero multiple of one row to another row}.
	
	\textbf{Example:} Let's say we wanted to add 4 times row 2 to row 3, i.e. we leave every row the same \textit{except} row 3, and we change row 3 by adding to it 4\text{R2} (shorthand: $\text{R3}=\text{R3}+4\text{R2}$).
	
	We could do this all at once, but to split it into steps, we could: 
	\begin{enumerate}[leftmargin=0.5in]
		\item compute $4\text{R2}$, i.e. $4\cdot\langle 1,0,0,-5\rangle=\langle 4,0,0,-20\rangle$;
		\item compute $\text{R3}$ plus $4\text{R2}$, i.e. $\langle 0,1,3,-1\rangle+\langle 4,0,0,-20\rangle=\langle 4,1,3,-21\rangle$; and
		\item form the new matrix having the same entries as $\sansM$ in rows 1, 2, and 4, and having $\text{R3}+4\text{R2}=\langle 4,1,3,-21\rangle$ as its third row. 
	\end{enumerate}

	Hence, the result is:
	
	\[
	\pmat{4 & 1 & -2 & 0 \\ 1 & 0 & 0 & -5 \\ 0 & 1 & 3 & -1 \\ 1 & 1 & -1 & -1} \xrightarrow{\text{R3}=\text{R3}+4\text{R2}}
	\pmat{4 & 1 & -2 & 0 \\ 1 & 0 & 0 & -5 \\ 4 & 1 & 3 & -21 \\ 1 & 1 & -1 & -1}
	\]
\end{enumerate}

\noindent Unsurprisingly, we can perform these three elementary row operations in succession to provide additional simplification. With a little foresight, this can yield a \textit{much simpler} matrix which is row-equivalent to the matrix we started with: \vspace{3mm}

\noindent\textbf{Example:}
\[
	\underbrace{\pmat{4 & 1 & -2 & 0 \\ 1 & 0 & 0 & -5 \\ 0 & 1 & 3 & -1 \\ 1 & 1 & -1 & -1}}_{\sansM}
	\xrightarrow{\text{R1}\leftrightarrow\text{R2}}
	\pmat{1 & 0 & 0 & -5 \\ 4 & 1 & -2 & 0 \\ 0 & 1 & 3 & -1 \\ 1 & 1 & -1 & -1}
	\xrightarrow{\text{R2}\leftrightarrow\text{R3}}
	\pmat{1 & 0 & 0 & -5 \\ 0 & 1 & 3 & -1 \\ 4 & 1 & -2 & 0 \\ 1 & 1 & -1 & -1}
	\xrightarrow{\text{R3}=\text{R3}-4\text{R1}}
	\pmat{1 & 0 & 0 & -5 \\ 0 & 1 & 3 & -1 \\ 0 & 1 & -2 & 20 \\ 1 & 1 & -1 & -1}
\]
\[
	\xrightarrow{\text{R3}=\text{R3}-\text{R2}}
	\pmat{1 & 0 & 0 & -5 \\ 0 & 1 & 3 & -1 \\ 0 & 0 & -5 & 21 \\ 1 & 1 & -1 & -1}
	\xrightarrow{\text{R3}=(-\frac{1}{5})\text{R3}}
	\pmat{1 & 0 & 0 & -5 \\ 0 & 1 & 3 & -1 \\ 0 & 0 & 1 & -\frac{21}{5} \\ 1 & 1 & -1 & -1}
	\longrightarrow\cdots
\]

\noindent Note that each of the above matrices is row-equivalent to $\sansM$.\vspace{3mm}

Moving forward, one of our main goals will be to perform these three elementary row operations in succession until we get to a matrix which is in \ul{Row Echelon Form (REF)} and/or \ul{Reduced Row Echelon Form (RREF)}.

\newpage

\subsection*{Row Echelon Form (REF)}
First, the definition:\vspace{3mm}

\noindent\textbf{Definition:} A matrix is in \ul{row echelon form (REF)} if it satisfies the following three properties:
\begin{enumerate}
	\item All nonzero rows are above any rows of all zeros.
	\item Each leading (nonzero) entry of a row is in a column to the right of the leading (nonzero) entry of the row above it.
	\item All entries in a column below a leading (nonzero) entry are zeros.
\end{enumerate}

As a remark, note that the entries \ul{above} the leading (nonzero) entries of such a matrix \ul{may or may not equal 0}. For instance, \textbf{both} of the following matrices are in REF:
\[
\sansA=\pmat{\fbox{1} & 1 & 1 & 2 \\ 0 & \fbox{2} & 3 & -1 \\ 0 & 0 & 0 & 0}\quad\quad\quad
\sansB=\pmat{\fbox{1} & 0 & 1 & 2 \\ 0 & \fbox{2} & 3 & -1 \\ 0 & 0 & 0 & 0}
\]

\noindent Note that an entry \fbox{in a box} denotes \ul{the leading nonzero entry of that particular row}.

\subsection*{Reduced Row Echelon Form (RREF)}
As it happens, \textit{neither} $\sansA$ nor $\sansB$ are in \ul{reduced row echelon form (RREF)}, because both have some mild simplifications that can be done to them.\vspace{3mm}

\noindent\textbf{Definition:} A matrix is in \ul{reduced row echelon form (RREF)} if it satisfies the following three properties:
\begin{enumerate}
	\item It is in REF;
	\item The leading (nonzero) entry in each row is 1.
	\item Each leading 1 is the \ul{only} nonzero entry in its column.
\end{enumerate}

In particular, we note that $\sansA$ fails both (2) and (3): The leading nonzero entry in row2 (i.e. the (2,2)-entry) is 2, and the column containing that 2 has other nonzero entries (namely, the 1 in the (1,2)-entry). However, we can put $\sansA$ in RREF using two elementary row operations:
\[
\sansA=\pmat{\fbox{1} & 1 & 1 & 2 \\ 0 & \fbox{2} & 3 & -1 \\ 0 & 0 & 0 & 0}
\xrightarrow{\text{R2}=(\frac{1}{2})\text{R2}}
\pmat{\fbox{1} & 1 & 1 & 2 \\ 0 & \fbox{1} & \frac{3}{2} & -\frac{1}{2} \\ 0 & 0 & 0 & 0}
\xrightarrow{\text{R1}=\text{R1}-\text{R2}}
\underbrace{\pmat{\fbox{1} & 0 & -\frac{1}{2} & \frac{5}{2} \\ 0 & \fbox{1} & \frac{3}{2} & -\frac{1}{2} \\ 0 & 0 & 0 & 0}}_{\text{RREF}}.
\]

\noindent For $\sansB$, one elementary row operation will get it into RREF:
\[
\sansB=\pmat{\fbox{1} & 0 & 1 & 2 \\ 0 & \fbox{2} & 3 & -1 \\ 0 & 0 & 0 & 0}
\xrightarrow{\text{R2}=(\frac{1}{2})\text{R2}}
\underbrace{\pmat{\fbox{1} & 0 & 1 & 2 \\ 0 & \fbox{1} & \frac{3}{2} & -\frac{1}{2} \\ 0 & 0 & 0 & 0}}_{\text{RREF}}
\]
\end{document}